% Options for packages loaded elsewhere
\PassOptionsToPackage{unicode}{hyperref}
\PassOptionsToPackage{hyphens}{url}
\documentclass[
  12pt,
]{article}
\usepackage{xcolor}
\usepackage[margin=1in]{geometry}
\usepackage{amsmath,amssymb}
\setcounter{secnumdepth}{-\maxdimen} % remove section numbering
\usepackage{iftex}
\ifPDFTeX
  \usepackage[T1]{fontenc}
  \usepackage[utf8]{inputenc}
  \usepackage{textcomp} % provide euro and other symbols
\else % if luatex or xetex
  \usepackage{unicode-math} % this also loads fontspec
  \defaultfontfeatures{Scale=MatchLowercase}
  \defaultfontfeatures[\rmfamily]{Ligatures=TeX,Scale=1}
\fi
\usepackage{lmodern}
\ifPDFTeX\else
  % xetex/luatex font selection
\fi
% Use upquote if available, for straight quotes in verbatim environments
\IfFileExists{upquote.sty}{\usepackage{upquote}}{}
\IfFileExists{microtype.sty}{% use microtype if available
  \usepackage[]{microtype}
  \UseMicrotypeSet[protrusion]{basicmath} % disable protrusion for tt fonts
}{}
\makeatletter
\@ifundefined{KOMAClassName}{% if non-KOMA class
  \IfFileExists{parskip.sty}{%
    \usepackage{parskip}
  }{% else
    \setlength{\parindent}{0pt}
    \setlength{\parskip}{6pt plus 2pt minus 1pt}}
}{% if KOMA class
  \KOMAoptions{parskip=half}}
\makeatother
\usepackage{graphicx}
\makeatletter
\newsavebox\pandoc@box
\newcommand*\pandocbounded[1]{% scales image to fit in text height/width
  \sbox\pandoc@box{#1}%
  \Gscale@div\@tempa{\textheight}{\dimexpr\ht\pandoc@box+\dp\pandoc@box\relax}%
  \Gscale@div\@tempb{\linewidth}{\wd\pandoc@box}%
  \ifdim\@tempb\p@<\@tempa\p@\let\@tempa\@tempb\fi% select the smaller of both
  \ifdim\@tempa\p@<\p@\scalebox{\@tempa}{\usebox\pandoc@box}%
  \else\usebox{\pandoc@box}%
  \fi%
}
% Set default figure placement to htbp
\def\fps@figure{htbp}
\makeatother
\setlength{\emergencystretch}{3em} % prevent overfull lines
\providecommand{\tightlist}{%
  \setlength{\itemsep}{0pt}\setlength{\parskip}{0pt}}
\usepackage{bookmark}
\IfFileExists{xurl.sty}{\usepackage{xurl}}{} % add URL line breaks if available
\urlstyle{same}
\hypersetup{
  hidelinks,
  pdfcreator={LaTeX via pandoc}}

\author{}
\date{\vspace{-2.5em}}

\begin{document}

\thispagestyle{empty}

\today

Editor\\
The R Journal\\
\bigskip

Dear Professor Tanaka, \bigskip

Please consider our article titled ``ShrinkageTrees: An R Package for
Bayesian Tree Ensembles for Survival Analysis and Causal Inference'' for
publication in the R Journal.

Bayesian tree ensembles for survival data are in high demand across
biostatistics, clinical research, and epidemiology, yet existing R
packages cover only part of the landscape. The \texttt{BART} package is
currently the only CRAN package that implements Bayesian tree ensembles
for survival outcomes, and it does not support interval-censored data,
causal forest models, or global--local shrinkage priors on the leaf
parameters. Our package \texttt{ShrinkageTrees} addresses all three
gaps. It extends BART to right-censored and interval-censored survival
outcomes via accelerated failure time models, provides the first
implementation of Bayesian causal forests for time-to-event endpoints in
R, and implements Horseshoe Forests --- tree ensembles with horseshoe
priors on the leaf step heights for adaptive shrinkage in
high-dimensional settings. The package also supports Dirichlet splitting
priors (DART) and allows both forms of regularisation to be combined
within a single model.

Beyond the methodology, we have invested in making
\texttt{ShrinkageTrees} a well-designed R package. It provides S3
classes with \texttt{print}, \texttt{summary}, \texttt{predict}, and
\texttt{plot} methods, integrates with \texttt{ggplot2} for
visualisation and \texttt{coda} for MCMC diagnostics, and supports
multi-chain parallel sampling out of the box. The accompanying paper is
dedicated to highlighting these software design choices and
demonstrating the package through a worked example on high-dimensional
ovarian cancer survival data.

We believe this article is a good fit for the R Journal because it
introduces a package that fills a clear gap in the R ecosystem for
Bayesian non-parametric survival analysis and causal inference, two
areas of active methodological and applied interest.

\bigskip
\bigskip

Regards,

Tijn Jacobs\\
Department of Mathematics\\
Vrije Universiteit Amsterdam\\
Amsterdam, The Netherlands\\
\href{mailto:t.jacobs@vu.nl}{\nolinkurl{t.jacobs@vu.nl}}

\bigskip

\end{document}
